% !Mode:: "TeX:UTF-8"

\chapter{实验与结果分析}

在第二章中，我们从算法与实现两个层面对基于 FSDP 通信钩子的
量化与稀疏化方法进行了系统梳理。然而，单纯从“比特数减少了多少”“稀
疏率达到多少”这样的静态指标，并不能回答通信压缩在实际训练中的关键
问题：在给定硬件与模型配置下，这些压缩方法究竟能在多大程度上降低通
信开销，又会在多大程度上伤害训练质量。一方面，压缩必然在梯度上引入
偏差和额外方差，可能导致损失曲线收敛变慢、最终困惑度（PPL）劣化；
另一方面，实现压缩与解压本身也需要额外计算与通信步骤，其带宽节省能
否真正体现为端到端吞吐（tokens/s）的提升，很大程度上取决于具体的通信
拓扑与算力–带宽比。

基于此，本章的目标可以概括为三个层面：（1）从训练质量角度，对比
baseline 无 压 缩 与 不 同 通 信 压 缩 算 法 在 相 同 训 练 步 数 或 epoch 内 的
Loss/step 收敛速度与平滑程度，以及在同源 Wikipedia held-out 验证集上的
最终 PPL 表现，对比各压缩算法在多大程度上牺牲了收敛质量；（2）从系
统效率角度，在完全对齐的模型、数据和优化配置下，比较各方法的吞吐
（tokens/s）、显存占用（allocated/reserved）以及各压缩算法的压缩比/稀疏
率，对比各压缩算法在缓解通信瓶颈、提升整体训练效率方面的表现；（3）
从机制解释角度，利用运行时记录的梯度压缩误差（相对 $\ell_2$ 误差
$\|\mathbf{g}-\hat{\mathbf{g}}\|_2/\|\mathbf{g}\|_2$）与上述质量、效率指标进行关联分析，体现不同算法在“压缩强度—
梯度扭曲程度—训练效果”三者之间形成的实际权衡。

\section{实验配置}

为保证对比的公平性，本章所有实验均在严格对齐的训练配置下进行：
保持训练数据、模型结构、优化器与学习率调度、随机种子、并行规模（worldsize）、
序列长度和 batch 配置完全一致，仅在通信相关参数如 comm\_hook
类型的选择以及各算法特有超参数上做差异化设置。所有结果均以无压缩的
FSDP baseline 作为参照系，重点报告相对于 baseline 的提升或下降幅度，
而非孤立的绝对数值，从而尽量减弱不同硬件环境和运行噪声对结论的干扰。
表~\ref{tab:exp-config}~说明了本章使用的核心实验配置。

\begin{table}[htbp]
\caption{压缩算法对比实验参数配置}
\label{tab:exp-config}
\vspace{0.5em}\centering\wuhao
\setlength{\tabcolsep}{3pt}
\renewcommand{\arraystretch}{1.5}
\begin{tabularx}{\textwidth}{
	>{\centering\arraybackslash}p{0.18\textwidth}
	>{\centering\arraybackslash}p{0.38\textwidth}
	>{\centering\arraybackslash}X
}
\toprule
参数项 & 配置值 & 说明 \\
\midrule
实验环境 & Ubuntu 24.04, 2$\times$A100 (40GB) & 操作系统与GPU硬件配置 \\
模型架构 & GPT-2 (Random Init) & 聚焦通信与计算特性 \\
训练精度 & FP32 & 全精度训练 \\
全局批大小 & 8 (Per-Device) & 单设备负载保持一致 \\
序列长度 & 1024 & 典型长序列训练场景 \\
显卡数量 & 2 & 单机2卡FSDP训练 \\
优化器 & AdamW + 线性衰减 & 各配置间保持一致 \\
训练步数 & 200 steps & 固定训练步数 \\
训练数据集 & Wikipedia En & Wikipedia词条语料 \\
验证数据集 & 同源Wikipedia held-out 5\% & 按样本划出约5\%作为验证集 \\
\bottomrule
\end{tabularx}
\renewcommand{\arraystretch}{1.0}
\setlength{\tabcolsep}{6pt}
\end{table}

通信压缩算法配置涵盖：无压缩baseline、int8 linear、Natural Compression、1-bit Seide（col\_size=256）、QSGD（$s$=4）、Threshold-v（自适应阈值，平均稀疏率$\le$0.01）、Sketched-SGD、Random-k（$\rho$=0.01）和Top-k（$\rho$=0.01）。

数据采集方面，我们对每次实验记录两类信息：一类是训练过程数据，
包括随 step 变化的训练损失 Loss/step、GPU 显存占用（allocated/reserved）
以及基于 wall‑clock 时间估算的吞吐（tokens/s）；另一类是通信与压缩误
差数据。前者用于刻画不同通信配置下的收敛行为和系统性能，后者则用于
量化压缩强度与训练效果之间的对应关系。

\section{实验结果}

在本节中，图~\ref{fig:loss-curves}~给出了不同通信压缩算法在表~\ref{tab:exp-config}~配置下的训练 loss
曲线，表~\ref{tab:exp-results}~则汇总了训练质量、系统吞吐与梯度压缩误差等关键指标。其
中 baseline 表示原始 FSDP 框架。

\begin{figure}[htbp]
\centering
\includegraphics[width=0.95\textwidth]{fig8}
\caption{不同通信压缩算法的训练损失曲线对比}
\label{fig:loss-curves}
\end{figure}

{
\wuhao
\setlength{\tabcolsep}{0pt}
\renewcommand{\arraystretch}{1.5}
\begin{longtable}{@{\extracolsep{\fill}}%
	>{\raggedright\arraybackslash}p{0.13\textwidth}
	>{\centering\arraybackslash}p{0.09\textwidth}
	>{\centering\arraybackslash}p{0.09\textwidth}
	>{\centering\arraybackslash}p{0.11\textwidth}
	>{\centering\arraybackslash}p{0.15\textwidth}
	>{\centering\arraybackslash}p{0.11\textwidth}
	>{\centering\arraybackslash}p{0.09\textwidth}
	>{\centering\arraybackslash}p{0.08\textwidth}
	>{\centering\arraybackslash}p{0.10\textwidth}@{}}
\caption{200 step 各通信压缩算法实验结果}
\label{tab:exp-results}\\
\toprule
\makecell[c]{压缩算法} & \makecell[c]{最后20\\个step\\的平均\\loss} & \makecell[c]{验证集\\Perplexity} & \makecell[c]{训练吞吐\\(token/s)} & \makecell[c]{压缩比/稀疏\\率} & \makecell[c]{平均梯\\度相对\\误差\\(L2范数)} & \makecell[c]{通信耗时\\(ms)} & \makecell[c]{计算-通\\信覆盖\\率} & \makecell[c]{每张\\GPU显\\存预留\\值} \\
\midrule
\endfirsthead

\multicolumn{9}{r}{表~\thetable（续表）}\\
\toprule
\makecell[c]{压缩算法} & \makecell[c]{最后20\\个step\\的平均\\loss} & \makecell[c]{验证集\\Perplexity} & \makecell[c]{训练吞吐\\(token/s)} & \makecell[c]{压缩比/稀疏\\率} & \makecell[c]{平均梯\\度相对\\误差\\(L2范数)} & \makecell[c]{通信耗时\\(ms)} & \makecell[c]{计算-通\\信覆盖\\率} & \makecell[c]{每张\\GPU显\\存预留\\值} \\
\midrule
\endhead

\midrule
\multicolumn{9}{r}{续下页}\\
\endfoot

\bottomrule
\endlastfoot

baseline & 6.9499 & 5324.56 & 10578.96 & 1x & - & 26564.78 & 0.604 & 37.57 \\
\makecell[l]{int8 linear} & 7.2162 & 5677.46 & 14321.25 & \makecell[c]{压缩比$\approx$4x} & 0.5922 & 13727.28 & 0.759 & 37.97 \\
\makecell[l]{natural\\compression} & 7.7415 & 6937.58 & 12936.38 & \makecell[c]{压缩比$\approx$3.56x} & 0.8967 & 17148.23 & 0.626 & 38.46 \\
\makecell[l]{1-bit\\seide} & 7.1214 & 5545.15 & 12384.80 & \makecell[c]{压缩比$\approx$32x} & 1.0642 & 11160.23 & 0.647 & 38.81 \\
QSGD & 7.7470 & 7037.46 & 11712.42 & \makecell[c]{压缩比$\approx$4x} & 1.1821 & 13630.69 & 0.755 & 38.54 \\
\makecell[l]{threshold v} & 7.3756 & 6186.17 & 12027.74 & \makecell[c]{平均稀疏率\\$\le$0.01} & 0.7164 & 12858.59 & 0.824 & 38.58 \\
sketch & 7.8575 & 7166.74 & 11994.57 & \makecell[c]{理论压缩\\比最高40x} & 0.9553 & 17291.79 & 0.590 & 38.63 \\
\makecell[l]{random k} & 8.0925 & 8053.88 & 13837.80 & \makecell[c]{稀疏率=0.01} & 7.2219 & 12306.74 & 0.814 & 38.64 \\
\makecell[l]{top k} & 7.3752 & 6071.02 & 14071.59 & \makecell[c]{稀疏率=0.01} & 0.8159 & 13343.68 & 0.805 & 37.85 \\
\end{longtable}
}

\subsection{训练质量、系统效率与误差关联分析}

从表 2-3 可以看出，在训练质量方面，baseline 的最后 20 个 step 平均
loss 为 6.95，作为参照；在所有压缩方法中，1-bit seide 的 7.12 最接近
baseline，其次是 int8 linear 的 7.22、threshold v 与 top k 的 7.38 左右，
验证集 PPL 也都维持在 4500--6200 的区间，说明在 4 $\times$ 量化或 1\% 稀疏
率下，只要方向信息保留得足够好，整体收敛质量仍可接受。natural
compression、QSGD 和 sketch 的 loss 分别升至 7.74、7.75 和 7.86，对应的
PPL 也更高，表明在当前配置下它们牺牲了一部分收敛精度来换取更高的压
缩比；而 random k 的平均 loss 达到 8.09、PPL 超过 8000，明显劣于其它
所有方法，反映出在仅保留 1\% 梯度且方向随机的设定下，噪声被大幅放大。

从系统效率与通信开销来看，baseline 的通信耗时约为 $2.66 \times 10^4$ ms，
对应吞吐 $1.06 \times 10^4$ tokens/s；所有压缩方法都在不同程度上缩短了通
信时间：1-bit seide 最激进，约 $1.12 \times 10^4$ ms，int8 linear、
random k 和 top k 也都将通信时间压缩到 $1.2 \times 10^4$--$1.4 \times 10^4$ ms
区间，并提高到了 $1.18 \times 10^4$--$1.40 \times 10^4$ tokens/s 的吞吐。其
中 top k 在保证较好收敛质量的同时达到最高吞吐。natural compression 与
QSGD 在通信时间和吞吐上略逊于 int8 linear，但仍明显优于 baseline；
threshold v 的通信时间介于 int8 和 top k 与 baseline 之间，在带宽节约与
loss 抬升之间取得中庸折中；sketch 虽然论文中标称最高可达 40$\times$ 总通
信压缩，但由于两轮通信与重建开销，其实际通信耗时在各方法中反而偏高。

在压缩误差与收敛质量的关系上，int8 linear、threshold v 与 top k 的平
均梯度相对 L2 误差均在 0.6--0.9 左右，loss 和 PPL 也处于相对温和的区
间；1-bit seide、QSGD、sketch 的误差接近或略高于 1，对应的 loss/PPL 进
一步抬升；而 random k 在 1\% 稀疏率下的平均梯度相对 L2 误差高达 7.22，
loss 与 PPL 均明显劣化，印证了对于激进且随机的稀疏方案在当前步数和
学习率下会显著拖慢收敛甚至破坏优化轨迹。至于 GPU 显存预留峰值，各
方法都聚集在 37--39 GB 的窄区间内，说明在本实验规模下，通信压缩对
FSDP 训练的显存占用几乎没有本质影响。

结合这些结果，在相近训练质量前提下，int8 量化、top-k 和 threshold-v
稀疏在“压缩比—通信时间—收敛质量”的折中上表现最为均衡，1-bit 与
sketch 适合“极限带宽受限、可接受更高损失”的场景，而 random k 则需要
更谨慎的学习率与训练步数调度才能发挥其高压缩率优势。

\section{本章小结}

本章在统一实验配置下，从训练质量、
系统效率和压缩误差三个维度对9种通信压缩算法进行了系统性对比评估。
实验结果表明：（1）int8量化、Top-k和Threshold-v稀疏化在综合权衡上表现最优；
（2）压缩误差与收敛质量之间存在明显的单调关系，
梯度L2误差控制在1.0以内时收敛质量可接受；（3）
通信压缩在双卡环境下可实现最高约35\%的吞吐提升，
所有方法对显存占用无本质影响。这些结论为实际部署通信压缩方案提供了决策依据。

